\begin{table}[tbp]
\centering
\small
\caption{Standing representation notes for the apparatus formalization.}
\label{tab:representation-notes}
\begin{tabularx}{\textwidth}{@{}>{\raggedright\arraybackslash}p{0.24\textwidth}
>{\raggedright\arraybackslash}p{0.28\textwidth}
>{\raggedright\arraybackslash}X@{}}
\toprule
Representation issue & Where it appears & Disclosure effect \\
\midrule
\(\vsrc\) structural encoding &
\(\Cref{sec:vsrc}\), App~\ref{app:formalization} &
The rank-\(2\) exterior relations are checked on the \(2^r\)-coordinate
model.  The paper claims calibration and consistency, not a full exterior
algebra or motivic descent formalization. \\
\addlinespace
\(\Q\) rank-\(1\) projector traces &
Finite-source and no-go statements &
Scalar-shadow residuals are finite rational projector computations; the
Lean result is the finite trace calculation, not an arithmetic
classification theorem. \\
\addlinespace
Identity-matrix action &
Schur-collapse statements &
Full-source adequacy reduces to \(I\)-action on an explicit finite carrier.
The external transport that admits the native source remains a carried
recognition obligation. \\
\bottomrule
\end{tabularx}
\end{table}
